% =====================================================================
%  Table 1 — A4 한 페이지 버전 (표를 페이지 중앙에 배치, 캡션은 A4 폭)
%  컴파일: tectonic -X compile table1_a4.tex   (또는 xelatex)
% =====================================================================
\documentclass[a4paper,11pt]{article}
\usepackage[margin=2.5cm]{geometry}
\usepackage{booktabs}
\usepackage{array}
\usepackage{siunitx}
\usepackage{kotex}
\usepackage{caption}
\captionsetup{labelfont=bf, font=small, skip=8pt}  % 캡션은 본문(A4) 폭으로 줄바꿈

\newcommand{\best}[1]{\textbf{#1}}
\pagestyle{empty}

\begin{document}
% --- 표를 페이지 세로 중앙에 배치 ---
\null\vfill
\begin{center}
\captionof{table}{폐사체 수거 작업의 단계별 성능 개선 결과. 초기 OpenVLA 모델은
폐사 개체를 탐지하고 접근할 수 있었으나 파지에 실패하였다. DAgger를 적용하여
파지 능력을 회복하였으며, 이후 객체 물리 특성을 개선함으로써 최종적으로 94\%의
작업 성공률을 달성하였다.}
\label{tab:headline}

\renewcommand{\arraystretch}{1.25}
\begin{tabular}{l l c c}
\toprule
Stage & Key change & Grasp & Full task \\
\midrule
Base                 & multi-view RGB, proprio       & 0/8          & 0\% \\
OFT                  & action chunking               & 0/8          & 0\% \\
Wrist depth          & wrist (z-depth)               & 0/8          & 0\% \\
\best{DAgger}        & on-policy relabeling          & \best{7/8}   & 0\% \\
\midrule
Firm-hold            & friction, inertia $\uparrow$  & 12/16        & 75\% \\
\best{Thick capsule} & grasp capsule $\uparrow$      & \best{17/18} & \best{94\%} \\
\bottomrule
\end{tabular}
\end{center}
\vfill\null

\end{document}
